\item La intensidad de una onda sonora a una distancia fija de una bocina que vibra a 1.00 kHz es de $0.600\text{ W/m}^2$.
\begin{enumerate}
    \item Determine la intensidad que resulta si la frecuencia se aumenta a 2.50 kHz mientras se mantiene una amplitud de desplazamiento constante.
    \item Calcule la intensidad si la frecuencia se reduce a 0.500 kHz y se duplica la amplitud de desplazamiento.
\end{enumerate}